\section{八、运行期硬件安全：TEE、内存加密与机密虚拟机}

安全启动只证明“被允许的软件获得了第一次执行权”。机器进入运行态后，处理器仍要隔离普通世界与安全世界、不同虚拟机、CPU 与 DMA 设备，并保护离开芯片封装的数据。运行期硬件安全因此不是一个单独开关，而是身份、地址转换、内存加密、完整性、设备访问和生命周期状态共同组成的通路。

\begin{pdfdiagram}[x=1cm,y=1cm]
  \node[hw state, minimum width=2.6cm] (normal) at (1.5,3.9) {普通世界\\进程/普通 VM};
  \node[hw control, minimum width=2.6cm] (secure) at (1.5,2.2) {安全世界/TEE\\受保护 VM};
  \node[hw compute, minimum width=3.0cm, minimum height=1.8cm] (cpu) at (5.2,3.05) {CPU 核与安全监控\\特权、页表、上下文 ID};
  \draw[hw bus] (normal) -- (cpu);
  \draw[hw bus] (secure) -- (cpu);

  \node[hw control, minimum width=3.0cm, minimum height=1.8cm] (memsec) at (8.8,3.05) {内存安全引擎\\加密、完整性、重放计数};
  \draw[hw bus] (cpu) -- (memsec);
  \node[hw memory, minimum width=2.4cm] (dram) at (12.0,3.9) {外部 DRAM\\密文与元数据};
  \node[hw control, minimum width=2.4cm] (iommu) at (12.0,2.2) {IOMMU\\设备域与权限};
  \draw[hw bus] (memsec) -- (dram);
  \draw[hw return] (iommu) -- (memsec);

  \node[hw block, minimum width=2.6cm] (dev) at (8.8,0.4) {DMA 设备\\队列与 PASID};
  \draw[hw bus] (dev) -- (11.0,0.4) -- (iommu.south);
  \node[hw label, text width=11.6cm] at (6.5,-0.8) {访问被接受必须同时满足：当前安全身份允许该页、IOMMU 允许该设备、密钥代际有效、完整性元数据匹配；其中任一条件失败都不能把明文或旧页面交给请求者。};
\end{pdfdiagram}
\diagramnote{运行期硬件安全的四本账。CPU 维护执行身份，页表和 IOMMU 维护地址权限，内存安全引擎维护密钥与完整性代际，设备请求携带来源身份。DRAM 中的密文只是其中一层，不能替代访问控制或完成顺序。}

\topic{TEE 用硬件状态区分两个执行世界}

\term{可信执行环境}{Trusted Execution Environment, TEE}为受保护代码建立独立的特权状态、内存属性和异常入口。以安全世界/普通世界结构为例，安全监控器在受控入口保存当前世界的寄存器与控制状态，切换页表、异常向量和访问权限，再进入另一世界。普通世界即使拥有操作系统最高权限，也不能直接映射标记为安全的物理区域或安全外设。

另一类 enclave 结构把保护范围缩小到进程中的一段代码和数据。处理器进入 enclave 时验证入口与线程状态，Cache line 离开片上保护边界前被加密并绑定地址或版本信息；异常发生时先保存受保护状态，再把控制权交给普通内核。TEE 保护的是特定威胁模型下的机密性与完整性，不自动证明 enclave 程序无漏洞，也不能消除时间、Cache、功耗和共享资源侧信道。

\topic{内存加密必须同时回答地址绑定和重放问题}

单纯用固定密钥加密 DRAM 会留下两个漏洞。第一，攻击者可把某地址的密文复制到另一地址；若加密未绑定物理地址，解密后可能仍得到合法但位置错误的数据。第二，攻击者可把旧密文写回原地址，形成重放。实际内存安全引擎通常把物理地址、地址空间或虚拟机密钥标识参与加密 tweak，并以计数器、Merkle tree 或其他完整性元数据验证版本。

完整性树的根必须位于片上可信状态或受保护的非易失状态。更新一条 Cache line 时，控制器不仅写密文，还要更新叶子计数器和通向根的认证路径；突然掉电可能使数据与元数据处于不同代际。因此机密内存若还承担持久化语义，必须规定原子更新、日志或恢复流程，不能把“已经加密”解释为“已经可靠持久化”。

\topic{机密虚拟机给每个来宾分配独立密钥与状态}

机密虚拟机技术把威胁模型扩展到不完全信任宿主机管理程序。硬件为每个来宾分配地址空间标识和内存加密密钥，二级页表仍负责地址转换，但宿主对来宾私有页的普通读取只得到密文或被拒绝。进入来宾时，处理器加载受保护控制结构；退出时，寄存器状态按规定保存，宿主只看到必要的退出原因和共享通信页。

来宾与宿主、设备或其他来宾交换数据时必须显式建立共享页。页面从私有转共享前，要停止相关 CPU 与设备访问、完成 Cache 与 TLB/IOTLB 失效、改变加密属性并清除旧内容；共享转私有时则要验证所有权、重新初始化并建立新密钥代际。只修改页表位而不关闭在途 DMA，会让旧设备翻译继续访问已经归属另一安全域的页框。

\topic{设备接入决定安全边界是否闭合}

设备若不能理解加密内存属性，驱动通常使用共享缓冲或 bounce buffer。CPU 把私有数据复制到共享页并发布描述符，设备 DMA 只访问共享页；完成后 CPU 验证长度与状态，再把结果复制回私有页。支持 PASID、ATS 或设备可信执行扩展的设备可以减少复制，但失效协议更复杂：CPU 页表、IOMMU、设备 TLB、内存密钥和队列代际必须一致推进。

\begin{longtable}{L{0.18\linewidth}L{0.31\linewidth}L{0.39\linewidth}}
\toprule
边界 & 必须建立的事实 & 常见错误 \\
\midrule
安全世界切换 & 寄存器、页表、异常入口与外设权限均属于目标世界 & 只切代码入口，遗留旧世界地址或中断状态 \\\tablerowrule
私有页映射 & 页归属、密钥标识与二级翻译一致 & 宿主重映射页框但密钥/版本未同步 \\\tablerowrule
DMA 共享 & IOMMU 仅允许共享页，描述符与长度已发布 & 设备使用旧 ATS 翻译写入已转私有页面 \\\tablerowrule
异常与调试 & 受保护寄存器不会经普通保存区或调试口泄露 & JTAG、崩溃转储或性能追踪越过生命周期策略 \\\tablerowrule
销毁与复用 & 密钥失效、Cache/TLB/IOTLB 清理、页面清零完成 & 旧密文或旧完成项被新安全域接受 \\
\bottomrule
\end{longtable}

\topic{密钥与调试生命周期从制造一直延伸到报废}

芯片可把根密钥来源放在熔丝、PUF 派生状态或受保护安全存储中。制造态允许测试和调试，部署态锁定量产接口，返修态是否重新开放必须由签名授权与物理策略共同决定。Boot ROM 读取生命周期状态，选择允许的启动密钥、调试能力和恢复入口；运行期安全监控器继续执行这些限制。

报废或安全域销毁时，最可靠的边界通常是使密钥不可恢复，而不是等待逐字节擦除大容量 DRAM 或存储介质。\term{密码学擦除}{cryptographic erase}先销毁数据加密密钥，再清理可访问元数据和共享缓冲；它能快速使旧密文失去解释能力，但仍要保证密钥没有副本停留在调试状态、崩溃转储或设备固件中。
